% Wspólny graf bazowy dla trzech przykładów.
\newcommand{\drawExampleGraph}{%
  \node[graph vertex] (a) at (0,1.8) {$v_1$};
  \node[graph vertex] (b) at (-1.1,0.7) {$v_2$};
  \node[graph vertex] (c) at (1.1,0.7) {$v_3$};
  \node[graph vertex] (d) at (-1.1,-0.7) {$v_4$};
  \node[graph vertex] (e) at (1.1,-0.7) {$v_5$};
  \draw[draw=black!30,line width=0.5pt] (a)--(b) (a)--(c) (b)--(c);
  \draw[draw=black!30,line width=0.5pt] (b)--(d) (c)--(d) (c)--(e) (d)--(e);
}

\begin{figure}[H]
  \centering
  \begin{subfigure}[t]{0.31\textwidth}
    \centering
    \begin{tikzpicture}[scale=0.82]
      \drawExampleGraph
      \draw[line width=1.4pt] (a)--(b)--(c)--(b)--(d);
    \end{tikzpicture}
    \caption{Droga (\textit{walk})\\[-1mm]
      $v_1,v_2,v_3,v_2,v_4$}
    \label{fig:przyklad-walk}
  \end{subfigure}
  \hfill
  \begin{subfigure}[t]{0.31\textwidth}
    \centering
    \begin{tikzpicture}[scale=0.82]
      \drawExampleGraph
      \draw[line width=1.4pt] (a)--(b)--(c)--(d)--(b);
    \end{tikzpicture}
    \caption{Droga prosta (\textit{trail})\\[-1mm]
      $v_1,v_2,v_3,v_4,v_2$}
    \label{fig:przyklad-trail}
  \end{subfigure}
  \hfill
  \begin{subfigure}[t]{0.31\textwidth}
    \centering
    \begin{tikzpicture}[scale=0.82]
      \drawExampleGraph
      \draw[line width=1.4pt] (a)--(c)--(e)--(d);
    \end{tikzpicture}
    \caption{Droga elementarna (\textit{path})\\[-1mm]
      $v_1,v_3,v_5,v_4$}
    \label{fig:przyklad-path}
  \end{subfigure}
  \caption{Przykłady ciągów wierzchołków w tym samym grafie}
  \label{fig:przyklady-walk-trail-path}
\end{figure}
